\appendix
\chapter{System Environment and Deployment Guide}

\section{Local Development Prerequisites}
The CodeQuest AI platform is engineered using modern, open-source web and game development technologies. To configure a complete local development node, the following dependencies must be provisioned:

\begin{itemize}
    \item \textbf{Operating System:} Windows 10/11 (64-bit), Ubuntu 22.04 LTS, or macOS 13+.
    \item \textbf{Game Engine:} Unity Editor Version 6000.0.x (Unity 6 LTS) with Universal Render Pipeline (URP) package.
    \item \textbf{Backend Runtime:} Python 3.11+ with \texttt{pip} package manager.
    \item \textbf{Container Runtime:} Docker Desktop 4.25+ or Docker Engine with Docker Compose v2.
    \item \textbf{Database:} PostgreSQL 15+ (provisioned via Docker container).
    \item \textbf{Identity Provider:} Keycloak 23.0+ running on port 8080 (optional for local guest testing).
\end{itemize}

\section{Service Orchestration via Docker Compose}
The backend infrastructure is orchestrated using Docker Compose. Listing~\ref{lst:compose_spec} illustrates the primary service definitions for the PostgreSQL database, FastAPI gateway, and isolated execution worker.

\begin{lstlisting}[language=bash, caption={Docker Compose Service Configuration}, label={lst:compose_spec}]
# Start core persistence and IAM services
docker compose -f docker-compose.infra.yml up -d postgres keycloak

# Initialize database schema migrations
alembic upgrade head

# Build the isolated Python evaluator container
docker build -t python-sandbox:latest ./PythonExecutionService/sandbox

# Launch backend API gateway and execution service
uvicorn app.main:app --host 0.0.0.0 --port 8000 --reload
\end{lstlisting}

\chapter{API Endpoint Specifications}

\section{Challenge and Submission Interfaces}
The CodeQuest AI application gateway exposes RESTful JSON endpoints documented under OpenAPI (Swagger) specifications. Table~\ref{tab:api_endpoints} outlines key API routes.

\begin{table}[htbp]
\centering
\caption{Key REST API Endpoints in CodeQuest AI Gateway}
\label{tab:api_endpoints}
\small
\begin{tabularx}{\textwidth}{lp{0.25\textwidth}p{0.2\textwidth}X}
\toprule
\textbf{Method} & \textbf{Endpoint Path} & \textbf{Auth Required} & \textbf{Description} \\
\midrule
\texttt{GET} & \path{/api/v1/health} & None & Returns service health status and database connectivity. \\
\midrule
\texttt{GET} & \path{/api/v1/challenges} & Bearer JWT / Guest & Lists available programming challenges categorized by difficulty. \\
\midrule
\texttt{GET} & \path{/api/v1/challenges/{id}} & Bearer JWT / Guest & Retrieves challenge prompt, starter code, and visible test cases. \\
\midrule
\texttt{POST} & \path{/api/v1/submissions} & Bearer JWT / Guest & Submits learner code for isolated evaluation and test verification. \\
\midrule
\texttt{GET} & \path{/api/v1/users/me/progress} & Bearer JWT & Retrieves completed challenge IDs, total score, and completion timestamps. \\
\midrule
\texttt{POST} & \path{/internal/evaluate} & Shared Secret Token & Inter-service endpoint invoked by gateway to dispatch Docker container run. \\
\bottomrule
\end{tabularx}
\end{table}

\section{Sample Submission Payload and Response}
Listing~\ref{lst:submission_json} provides a representative HTTP request and response interaction between the Unity client HUD and the FastAPI backend.

\begin{lstlisting}[language=bash, caption={Sample Code Submission Request and Evaluator Response Payload}, label={lst:submission_json}]
# POST /api/v1/submissions
{
  "challenge_id": "ch_var_01",
  "language": "python",
  "code": "player_name = 'Hero'\nlevel = 1\nprint(f'{player_name} is level {level}')"
}

# HTTP 200 OK Response
{
  "submission_id": "sub_9841f32a",
  "status": "PASSED",
  "tests_total": 3,
  "tests_passed": 3,
  "execution_time_ms": 342,
  "stdout": "Hero is level 1\n",
  "stderr": "",
  "feedback": "Outstanding! All variable assignments and print assertions verified successfully."
}
\end{lstlisting}
